\documentclass[conference]{IEEEtran}
\IEEEoverridecommandlockouts

\usepackage{cite}
\usepackage{amsmath,amssymb,amsfonts}
\usepackage{graphicx}
\usepackage{xcolor}
\usepackage{booktabs}
\usepackage{tikz}
\usepackage{microtype}
\usepackage{hyperref}
\hypersetup{colorlinks=true,linkcolor=blue,citecolor=blue,urlcolor=blue}

\setlength{\parskip}{0pt}
\setlength{\floatsep}{4pt}
\setlength{\textfloatsep}{4pt}
\setlength{\intextsep}{4pt}
\setlength{\abovecaptionskip}{2pt}
\setlength{\belowcaptionskip}{0pt}

\begin{document}

\title{Nexus: A Multimodal Spatiotemporal Retrieval-Augmented\\
Generation Engine for Educational Video Understanding}

\author{
    \IEEEauthorblockN{Kavya Gujral}
    \IEEEauthorblockA{
        \textit{Department of Computer Science and Engineering}\\
        \textit{[University Name]} \quad [email@university.edu]
    }
}

\maketitle

% ── Abstract ──────────────────────────────────────────────────────────────────
\begin{abstract}
Recorded lectures are notoriously hard to search. You can ctrl+F a PDF, but
you cannot ctrl+F a video. \textbf{Nexus} is our attempt to fix that: a
multimodal RAG engine that lets students ask natural language questions about
educational recordings and get back exact timestamps, frame screenshots, and
synthesised answers. Under the hood, the system fuses speech (OpenAI Whisper
plus a PyTorch GRU for sequential context), vision (a pretrained ResNet-50
CNN), and temporal position (sinusoidal encodings) into 512-dimensional
vectors --- all joined through a single-head cross-modal attention mechanism
written from scratch in NumPy, without any deep learning framework. These
vectors live in a ChromaDB store, and queries come in as plain English.
\end{abstract}

\begin{IEEEkeywords}
RAG, multimodal learning, cross-modal attention, spatiotemporal grounding,
ChromaDB, Whisper, ResNet-50, GRU, educational video retrieval
\end{IEEEkeywords}

% ── 1. Introduction ───────────────────────────────────────────────────────────
\section{Introduction}

Consider a two-hour Zoom recording of a software architecture lecture. A
student trying to revisit the part about microservice communication has,
realistically, two options: scrub through the video manually, or hope the
transcript search surfaces the right segment. Neither works well. Transcript
search ignores what was \emph{on screen} --- and in technical lectures, the
slide or diagram often carries more information than the narration.

Nexus was built around a specific intuition: a lecture moment is only fully
described by three things together --- what was \emph{said}, what was
\emph{shown}, and \emph{when} it happened. Strip any of those out and you
lose fidelity. So rather than building yet another transcript search tool,
we encode all three into a joint embedding and let cosine similarity do the
heavy lifting.

The contributions are four-fold. First, a pure-NumPy single-head cross-modal
attention (CLO01/CLO02) --- no PyTorch, no autograd, every matrix multiply
visible and testable. Second, a two-second window alignment scheme that
handles the messy reality of Whisper segments that don't align neatly to
frame boundaries. Third, an eight-stage asynchronous ML pipeline served
through FastAPI with real-time frontend progress reporting. And fourth, a
dual-database design that separates relational metadata (PostgreSQL) from
vector search (ChromaDB).

% ── 2. System Architecture ────────────────────────────────────────────────────
\section{System Architecture}

\begin{figure}[ht]
\centering
\begin{tikzpicture}[
    box/.style={rectangle, draw, rounded corners=2pt, minimum width=2.55cm,
                minimum height=0.42cm, font=\scriptsize, fill=blue!8, align=center},
    sbox/.style={rectangle, draw, rounded corners=2pt, minimum width=1.2cm,
                minimum height=0.42cm, font=\scriptsize, fill=green!8, align=center},
    arrow/.style={->, >=stealth, thin},
    node distance=0.55cm
]
\node[box] (vid)    {Input Video};
\node[box, below left=0.45cm and 0.3cm of vid]  (fr)  {OpenCV Frames\\$\Delta t{=}2$\,s};
\node[box, below right=0.45cm and 0.3cm of vid] (au)  {FFmpeg Audio\\16\,kHz};
\node[box, below=0.45cm of fr]  (rn)  {ResNet-50 CNN\\512-dim};
\node[box, below=0.45cm of au]  (wh)  {Whisper STT\\+ GRU};
\node[box, below=0.9cm of vid]  (al)  {2\,s Window Alignment};
\node[box, below=0.45cm of al]  (at)  {NumPy Cross-Modal Attn};
\node[box, below=0.45cm of at]  (ch)  {ChromaDB 512-dim};
\node[sbox, below left=0.45cm and 0.5cm of ch]  (q)  {NL Query};
\node[sbox, below right=0.45cm and 0.5cm of ch] (r)  {Answer +\\Timestamp};

\draw[arrow] (vid) -- (fr);
\draw[arrow] (vid) -- (au);
\draw[arrow] (fr)  -- (rn);
\draw[arrow] (au)  -- (wh);
\draw[arrow] (rn)  -- (al);
\draw[arrow] (wh)  -- (al);
\draw[arrow] (al)  -- (at);
\draw[arrow] (at)  -- (ch);
\draw[arrow] (q)   -- (ch);
\draw[arrow] (ch)  -- (r);
\end{tikzpicture}
\vspace{-2pt}
\caption{Nexus end-to-end spatiotemporal RAG pipeline.}
\label{fig:arch}
\end{figure}

Fig.~\ref{fig:arch} shows the full flow. On the visual side, OpenCV samples
a frame every two seconds exactly --- $t \in \{0.0, 2.0, 4.0, \ldots\}$ ---
and each frame goes through a pretrained \textbf{ResNet-50} backbone. The
2048-dim penultimate features get projected down to $\mathbb{R}^{512}$, and
a lightweight classification head simultaneously guesses one of five visual
categories: \textit{architecture diagram}, \textit{code snippet},
\textit{slide}, \textit{whiteboard}, or \textit{talking-head} (CLO04).

The audio side is a bit more involved. FFmpeg strips the audio track to
16\,kHz mono, \textbf{Whisper} (\texttt{base} model) transcribes it into
timestamped segments, and \texttt{all-MiniLM-L6-v2} embeds each segment's
text into 512 dimensions. Those embeddings then run sequentially through
a bidirectional \textbf{PyTorch GRU} with hidden dimension 256 --- because
without some sequential modelling, the representation for ``as we mentioned
earlier'' has no idea what was mentioned earlier (CLO05):
\begin{equation}
h_t = \text{GRU}(x_t,\, h_{t-1}).
\end{equation}

% ── 3. Core Technical Contributions ──────────────────────────────────────────
\section{Core Technical Contributions}

\subsection{Two-Second Temporal Window Alignment}

Here is the practical problem: a Whisper segment might run from $t=1.3$\,s
to $t=5.1$\,s, but our frame grid runs on clean two-second boundaries.
We define windows $W_k = [2k,\, 2k{+}2)$ and assign speech to windows
proportionally by overlap:
\begin{equation}
  \text{Overlap}(W_k, S_j)=\max\!\bigl(0,\,
  \min(t_{k,e},\,t_e^j)-\max(t_{k,s},\,t_s^j)\bigr).
\end{equation}
Each window's audio vector is then a weighted average of all overlapping
segment embeddings. It is not a particularly glamorous component, but
getting this wrong corrupts every downstream fusion result --- so we
wrote dedicated unit tests for it.

\subsection{Pure NumPy Cross-Modal Attention (CLO01/CLO02)}

The fusion itself is a single-head scaled dot-product attention where
the GRU-encoded audio acts as the query and the ResNet visual features
supply the keys and values. Everything is in NumPy:
\begin{align}
  Q &= A W_Q,\quad K = V W_K,\quad V_l = V W_V \\
  S &= \tfrac{QK^\top}{\sqrt{d_k}},\quad
  \alpha = \tfrac{\exp(S-\max S)}{\textstyle\sum_j \exp(S_j-\max S)} \\
  Z &= \alpha V_l + A,\quad
  \hat{Z} = \tfrac{Z-\mu}{\sqrt{\sigma^2+\varepsilon}}\odot\gamma+\beta.
\end{align}

Weight matrices $W_Q, W_K, W_V \in \mathbb{R}^{512\times 512}$ use Xavier
uniform initialisation \cite{glorot2010init}. The residual addition $Z = \alpha V_l + A$
matters --- without it, the audio signal gets washed out entirely when
visual attention weights are dominant. Temporal position is injected via
sinusoidal encodings before fusion:
\begin{equation}
  \text{PE}(t,2i)=\sin\!\Bigl(\tfrac{t}{10000^{2i/d}}\Bigr),\quad
  \text{PE}(t,2i{+}1)=\cos\!\Bigl(\tfrac{t}{10000^{2i/d}}\Bigr).
\end{equation}
The fused vector is $\ell_2$-normalised to unit length before writing to
ChromaDB.

\subsection{Retrieval and Answer Synthesis}

Query embeddings use the same SentenceTransformer. ChromaDB returns the
top-$k$ results by cosine distance; anything below $\tau_{\min}=0.12$
similarity is discarded --- the system says ``nothing relevant found''
rather than confabulating an answer. Retrieved windows within six seconds
of each other get merged into a single \texttt{TopicSegment}, duplicate
sentences across adjacent windows are deduplicated, and the result is
formatted with a prominently displayed timestamp anchor.

% ── 4. Implementation ────────────────────────────────────────────────────────
\section{Implementation}

\begin{table}[ht]
\centering
\caption{Technology stack}
\label{tab:stack}
\setlength{\tabcolsep}{4pt}
\begin{tabular}{@{}lll@{}}
\toprule
\textbf{Layer} & \textbf{Technology} & \textbf{Role} \\ \midrule
Backend API        & FastAPI (Python 3.11)  & Async REST + pipeline jobs \\
Speech-to-Text     & Whisper \texttt{base}  & Timestamped segments \\
Visual Encoder     & ResNet-50 (pretrained) & Frame features + scene label \\
Sequential Context & PyTorch GRU            & Discourse-aware audio vectors \\
Attention          & Pure NumPy             & Cross-modal fusion \\
Vector DB          & ChromaDB (persistent)  & 512-dim cosine retrieval \\
Text Embeddings    & all-MiniLM-L6-v2       & Query + transcript embed \\
Relational DB      & PostgreSQL + asyncpg   & Video, job, status metadata \\
Frontend           & React 19 + Vite + TS   & Upload, pipeline UI, query \\
\bottomrule
\end{tabular}
\end{table}

Table~\ref{tab:stack} gives the full stack. The backend spans 15 Python
modules; the frontend has components for drag-and-drop upload, a live
eight-stage pipeline progress tracker, a query bar with preset chips, and
multimodal result cards that show the frame thumbnail alongside the
synthesised answer and a click-to-seek timestamp.

One design detail worth mentioning: Whisper intentionally runs on CPU
even on MPS-equipped Apple Silicon machines, because fp16 inference on
Apple's Neural Engine produces unreliable results with the current
\texttt{openai-whisper} release.

% ── 5. Evaluation ────────────────────────────────────────────────────────────
\section{Evaluation}

Retrieval quality is measured with Recall@$k$ ($k \in \{1,3,5\}$) using a
$\tau=3$\,s timestamp tolerance, and with MRR:
$\text{MRR}=\frac{1}{|Q|}\sum_q \frac{1}{\text{rank}_q}$.
Visual classification uses Top-1 accuracy and Macro-F1 across the five
educational scene categories.

\begin{table}[ht]
\centering
\caption{Mathematical unit test coverage (\texttt{pytest})}
\label{tab:tests}
\setlength{\tabcolsep}{3pt}
\begin{tabular}{@{}ll@{}}
\toprule
\textbf{Test} & \textbf{Property Verified} \\ \midrule
Softmax stability     & $\sum_j\alpha_j=1$; safe with $S_i\in[-500,1000]$ \\
LayerNorm             & $\mu\approx 0$, $\sigma^2\approx 1$ along feature axis \\
Single-head attention & Output shape; single-item $\alpha=1.0$ \\
Masked attention      & Masked weight $<10^{-6}$; unmasked $\approx 1.0$ \\
Temporal PE           & $\text{PE}\in[-1,1]$, shape $(T,d)$ correct \\
End-to-end fusion     & $\|\hat{Z}\|_2=1.0$ \\
Window alignment      & Overlap arithmetic, boundary edge cases \\
ChromaDB retrieval    & Insert, search, Recall@k pipeline \\
\bottomrule
\end{tabular}
\end{table}

Every mathematical claim in Section~III has a corresponding test in
Table~\ref{tab:tests}. That said, formal benchmark evaluation against a
labelled query set is still in progress --- the evaluation module in
\texttt{backend/evaluation/metrics.py} returns an explicit
\textit{``Evaluation dataset not yet provided''} status rather than
inventing numbers.

% ── 6. Results & Discussion ───────────────────────────────────────────────────
\section{Results \& Discussion}

At the time of writing, both servers are running: FastAPI at
\texttt{localhost:8000} and the Vite frontend at \texttt{localhost:5173}.
ChromaDB holds 284 indexed spatiotemporal chunks from ingested test footage,
and qualitative tests show that the similarity gate ($\tau_{\min}=0.12$)
correctly rejects off-topic queries without returning garbage results.

We deliberately chose NumPy over a PyTorch \texttt{nn.Module} for the
attention layer. Yes, it is slower and slightly more painful to write, but
it means every intermediate value --- the raw score matrix, the softmax
weights, the residual sum --- is a plain numpy array you can inspect in a
debugger or test against an analytical expectation. That transparency was
more valuable for this project than the speed gain.

\noindent\textbf{Known limitations.} PostgreSQL connectivity requires the
\texttt{greenlet} package (the system falls back to an in-memory store
without it). The ResNet classifier is tuned for educational content and
likely won't generalise well elsewhere without fine-tuning.

\noindent\textbf{Planned next steps.} A multi-head attention variant
is already stubbed in \texttt{numpy\_attention.py} as optional future
work. We also plan to benchmark against a labelled query set and explore
speaker diarisation for multi-participant recordings.

% ── 7. Conclusion ─────────────────────────────────────────────────────────────
\section{Conclusion}
Nexus is a working multimodal RAG system that treats educational video as
a first-class retrieval target. The combination of Whisper, ResNet-50, a
bidirectional GRU, and a hand-rolled NumPy attention layer produces
512-dimensional embeddings that support meaningful semantic search over
lecture content. The groundwork --- pipeline, vector store, evaluation
harness, unit tests --- is in place; quantitative benchmarking is the
obvious next step.

% ── References ────────────────────────────────────────────────────────────────
\begin{thebibliography}{00}
\bibitem{lewis2020rag} P.~Lewis et al., ``Retrieval-augmented generation for knowledge-intensive NLP tasks,'' \textit{NeurIPS}, vol.~33, pp.~9459--9474, 2020.
\bibitem{vaswani2017attention} A.~Vaswani et al., ``Attention is all you need,'' \textit{NeurIPS}, vol.~30, 2017.
\bibitem{radford2023whisper} A.~Radford et al., ``Robust speech recognition via large-scale weak supervision,'' \textit{ICML}, pp.~28492--28518, 2023.
\bibitem{he2016resnet} K.~He et al., ``Deep residual learning for image recognition,'' \textit{CVPR}, pp.~770--778, 2016.
\bibitem{cho2014gru} K.~Cho et al., ``Learning phrase representations using RNN encoder-decoder,'' \textit{EMNLP}, pp.~1724--1734, 2014.
\bibitem{glorot2010init} X.~Glorot and Y.~Bengio, ``Understanding the difficulty of training deep feedforward neural networks,'' \textit{AISTATS}, pp.~249--256, 2010.
\end{thebibliography}

\end{document}
